\appendix
\begin{frame}[plain,noframenumbering]
    \centering
    \vspace{2cm}
    {\Huge\textbf{Anhang}}
\end{frame}

\begin{frame}{Eigenschaften von Markov-Ketten}
    \begin{columns}
        \column{0.5\textwidth}
        \begin{itemize}[<+->]
            \item \textbf{Irreduzibilität:} Von jedem Zustand kann man jeden anderen erreichen
            \item \textbf{Aperiodizität:} Rückkehr zu einem Zustand nicht nur in festen Zeitabständen
            \item \textbf{Absorbierende Zustände:} Zustände, die nicht mehr verlassen werden
        \end{itemize}

        \column{0.5\textwidth}
        \begin{itemize}[<+->]
            \item \textbf{Rekurrenz:} Wahrscheinlichkeit, wieder in einen Zustand zurückzukehren
            \item \textbf{Transienz:} Zustände, die irgendwann nicht mehr erreicht werden
            \item \textbf{Ergodizität:} Positive Rekurrenz und Aperiodizität
        \end{itemize}
    \end{columns}
\end{frame}

\begin{frame}{Absorbierende Zustände}
    \begin{columns}
        \column{0.6\textwidth}
        Beispiel:
        \[
            P = \begin{pmatrix}
                1   & 0   \\
                0.3 & 0.7 \\
            \end{pmatrix}
        \]

        \column{0.4\textwidth}
        \begin{figure}[H]
            \centering
            \begin{tikzpicture}[node distance=2cm, auto, thick, scale=0.7, transform shape]
                \node[state] (1) {1};
                \node[state] (2) [below of=1] {2};

                \path[->]
                (1) edge[loop right] node {1.0} (1)
                (2) edge[bend left] node {0.3} (1)
                (2) edge[loop right] node {0.7} (2);
            \end{tikzpicture}
        \end{figure}
    \end{columns}

    \only<2->{
        \begin{itemize}
            \item Zustand 1 ist \textbf{absorbierend}: $P_{1,1} = 1$
            \item Wenn die Kette in Zustand 1 ist, bleibt sie dort für immer
            \item Zustand 2 ist \textbf{transient}: Irgendwann wird er verlassen und nie wieder erreicht
        \end{itemize}
    }
\end{frame}

\begin{frame}{Grenzen von Markov-Ketten}
    \begin{itemize}[<+->]
        \item \textbf{Markov-Eigenschaft:} Nicht immer realistisch, manchmal hängt die Zukunft von mehr als nur dem aktuellen Zustand ab
        \item \textbf{Feste Übergangswahrscheinlichkeiten:} In der Realität können sich Wahrscheinlichkeiten ändern
        \item \textbf{Komplexe Abhängigkeiten:} Können nicht abgebildet werden
        \item \textbf{Zustandsraum:} Bei vielen Zuständen wird die Berechnung aufwändig
        \item \textbf{Keine Kausalität:} Beschreiben nur Wahrscheinlichkeiten, keine Ursachen
    \end{itemize}
\end{frame}

\begin{frame}{App-Nutzung: Stationäre Verteilung}
    Für die App-Nutzung: $P = \begin{pmatrix} 0.7 & 0.3 \\ 0.3 & 0.7 \end{pmatrix}$

    \begin{enumerate}[<+->]
        \item Aus $\pi_1 = 0.7\pi_1 + 0.3\pi_2$ folgt $0.3\pi_1 = 0.3\pi_2$ → $\pi_1 = \pi_2$
        \item Mit $\pi_1 + \pi_2 = 1$ folgt $\pi_1 = \pi_2 = 0.5$
    \end{enumerate}

    \only<3->{
        Langfristig werden beide Apps gleich häufig genutzt.
    }
\end{frame}


\begin{frame}{Anwendungen von Markov-Ketten}
    \begin{columns}
        \column{0.5\textwidth}
        \begin{figure}[H]
            \centering
            \includegraphics[width=0.25\pagewidth]{"Grundlagen_Info/10_Aus_Daten_lernen/Figures/intro_chatGPT"}
        \end{figure}
        \begin{itemize}[<+->]
            \item \textbf{Textgenerierung} in ChatGPT und anderen KI-Modellen
            \item \textbf{Wettervorhersage}
            \item \textbf{PageRank-Algorithmus} von Google
            \item \textbf{Biologische Systeme} (Gensequenzen)
        \end{itemize}

        \column{0.5\textwidth}
        \begin{itemize}[<+->]
            \item \textbf{Finanzmodelle} (Aktienkurse)
            \item \textbf{Verkehrssimulationen}
            \item \textbf{Warteschlangensysteme}
            \item \textbf{Spieltheorie} (Entscheidungsprozesse)
            \item \ldots
        \end{itemize}
    \end{columns}
\end{frame}


\begin{frame}{Textvorhersage mit Markov-Ketten}
    \begin{itemize}[<+->]
        \item Markov-Ketten können genutzt werden, um \textbf{Texte zu generieren}.
        \item Die Idee: Jedes Wort (oder Zeichen) ist ein Zustand.
        \item Die Wahrscheinlichkeit für das nächste Wort hängt nur vom aktuellen Wort ab.
    \end{itemize}
    \vspace{0.5cm}
    \begin{block}{Beispiel:}
        \textit{„Ich mag Pizza und ich mag Pasta.“}
    \end{block}
    \vspace{0.3cm}
    \only<3->{
        \begin{itemize}
            \item Übergänge (vereinfacht):\\
                  \texttt{Ich} $\rightarrow$ \texttt{mag} $\rightarrow$ \texttt{Pizza} $\rightarrow$ \texttt{und} $\rightarrow$ \texttt{ich} $\rightarrow$ \texttt{mag} $\rightarrow$ \texttt{Pasta}
        \end{itemize}
    }
\end{frame}

\begin{frame}{Visualisierung: Markov-Kette für Text}
    \centering
    \begin{tikzpicture}[node distance=2.5cm, auto, thick, >=stealth']
        \node[state] (Ich) {\texttt{Ich}};
        \node[state] (mag) [right of=Ich] {\texttt{mag}};
        \node[state] (Pizza) [right of=mag] {\texttt{Pizza}};
        \node[state] (und) [below of=Pizza] {\texttt{und}};
        \node[state] (Pasta) [above of=Pizza] {\texttt{Pasta}};

        \path[->]
        (Ich) edge node {1.0} (mag)
        (mag) edge node[below] {0.5} (Pizza)
        (mag) edge node[below] {0.5} (Pasta)
        (Pizza) edge node {1.0} (und)
        (und) edge node {1.0} (Ich);
    \end{tikzpicture}
    \vspace{0.5cm}

    \only<2->{
        \begin{itemize}
            \item Jeder Pfeil steht für eine mögliche Fortsetzung im Text.
            \item Die Wahrscheinlichkeiten werden aus Beispielsätzen gezählt.
        \end{itemize}
    }
\end{frame}

\begin{frame}[fragile]{Kurze Textgenerierung (Beispiel)}
    \begin{itemize}[<+->]
        \item Start: \texttt{Ich}
        \item Nächstes Wort: \texttt{mag}
        \item Danach: zufällig \texttt{Pizza} oder \texttt{Pasta}
        \item usw.
    \end{itemize}
    \vspace{0.3cm}
    \begin{block}{Mögliche generierte Sätze:}
        \texttt{Ich mag Pizza und ich mag Pasta.}\\
        \texttt{Ich mag Pasta.}
    \end{block}
    \vspace{0.3cm}
    \only<5->{
        Markov-Ketten können einfache, aber oft überraschende Texte erzeugen!
    }
\end{frame}

\begin{frame}
    \begin{itemize}
        \item \aufgabebox{Text-Generierung mit Markov-Ketten}
        \item \challengebox{1.3}
    \end{itemize}
\end{frame}
\begin{frame}{Lernkontrolle: Markov-Ketten für Text}
    \begin{myexercise}

        \texttt{Hallo Welt Info Hallo Info Welt Hallo Hallo Welt Info Info Welt Hallo}

        \begin{enumerate}
            \item Leiten Sie aus dem Text eine Übergangsmatrix für die Wörter \texttt{Hallo}, \texttt{Welt} und \texttt{Info} her.
            \item Was ist der wahrscheinlichste Text mit 3 Wörtern, wenn Sie mit \texttt{Hallo} starten?
        \end{enumerate}
        \textit{Hinweis: Überlegen Sie sich sinnvolle Übergangswahrscheinlichkeiten!}

    \end{myexercise}
\end{frame}

\begin{frame}
    \begin{enumerate}
        \item \textbf{Übergangszählung:}
              \begin{table}[H]
                  \centering
                  \begin{tabular}{c|ccc}
                                          & \texttt{Hallo} & \texttt{Welt} & \texttt{Info} \\
                      \hline
                      \texttt{Hallo} (5x) & 1              & 2             & 1             \\
                      \texttt{Welt} (4x)  & 2              & 0             & 2             \\
                      \texttt{Info} (4x)  & 1              & 2             & 1             \\
                  \end{tabular}
              \end{table}

              \textbf{Übergangsmatrix:}

              Für jede Zeile normieren wir auf die Summe der Übergänge von diesem Wort:

              \[
                  P = \begin{pmatrix}
                                     & \texttt{Hallo} & \texttt{Welt} & \texttt{Info} \\
                      \texttt{Hallo} & \frac{1}{4}    & \frac{2}{4}   & \frac{1}{4}   \\
                      \texttt{Welt}  & \frac{2}{4}    & 0             & \frac{2}{4}   \\
                      \texttt{Info}  & \frac{1}{4}    & \frac{2}{4}   & \frac{1}{4}   \\
                  \end{pmatrix}
              \]

        \item \textbf{Wahrscheinlichster Text mit 3 Wörtern, Start \texttt{Hallo}:}

              Von \texttt{Hallo} aus ist der wahrscheinlichste Übergang zu \texttt{Welt} (2 von 4 Übergängen), danach von \texttt{Welt} zu \texttt{Hallo} oder \texttt{Info} (je 2 von 4 Übergängen).

              \begin{itemize}
                  \item \texttt{Hallo} $\rightarrow$ \texttt{Welt} ($0.5$), dann \texttt{Welt} $\rightarrow$ \texttt{Hallo} ($0.5$):\\
                        \texttt{Hallo Welt Hallo}
                  \item \texttt{Hallo} $\rightarrow$ \texttt{Welt} ($0.5$), dann \texttt{Welt} $\rightarrow$ \texttt{Info} ($0.5$):\\
                        \texttt{Hallo Welt Info}
              \end{itemize}

              Beide sind gleich wahrscheinlich und am wahrscheinlichsten.

              \textbf{Übergangsmatrix:}

              Für jede Zeile normieren wir auf die Summe der Übergänge von diesem Wort:

              \[
                  P = \begin{pmatrix}
                                     & \texttt{Hallo} & \texttt{Welt} & \texttt{Info} \\
                      \texttt{Hallo} & 0              & \frac{3}{6}   & \frac{3}{6}   \\
                      \texttt{Welt}  & \frac{2}{5}    & 0             & \frac{3}{5}   \\
                      \texttt{Info}  & \frac{3}{5}    & \frac{2}{5}   & 0             \\
                  \end{pmatrix}
              \]

        \item \textbf{Wahrscheinlichster Text mit 3 Wörtern, Start \texttt{Hallo}:}

              Von \texttt{Hallo} aus gibt es gleichwahrscheinliche Übergänge zu \texttt{Welt} und \texttt{Info}. Danach ist der wahrscheinlichste nächste Schritt jeweils der Übergang mit der höchsten Wahrscheinlichkeit:

              \begin{itemize}
                  \item \texttt{Hallo} $\rightarrow$ \texttt{Welt} ($0.5$), dann \texttt{Welt} $\rightarrow$ \texttt{Info} ($0.6$):\\
                        \texttt{Hallo Welt Info}
                  \item \texttt{Hallo} $\rightarrow$ \texttt{Info} ($0.5$), dann \texttt{Info} $\rightarrow$ \texttt{Hallo} ($0.6$):\\
                        \texttt{Hallo Info Hallo}
              \end{itemize}

              Beide sind gleich wahrscheinlich und am wahrscheinlichsten.
    \end{enumerate}
\end{frame}

\begin{frame}{Visualisierung von Markov-Ketten}
    Markov-Ketten können als \textbf{gerichtete Graphen} dargestellt werden:

    \begin{columns}

        % \begin{table}[H]
        %     \centering
        %     \begin{tabular}{c c|c|c}
        %                                                       &                  & \multicolumn{2}{c}{\textbf{Nach}}                    \\
        %                                                       &                  & \emoji{sun}                       & \emoji{umbrella} \\\midrule\midrule
        %         \multirow{2}{*}{\rotatebox{90}{\textbf{Von}}} & \emoji{sun}      & 0.8                               & 0.2              \\\cmidrule{2-4}
        %                                                       & \emoji{umbrella} & 0.4                               & 0.6              \\
        %     \end{tabular}
        %     \caption*{Übergangsmatrix}
        % \end{table}
        % \column{.05\textwidth}
        % \only<2->{
        %     $\Rightarrow$
        % }
        \column{.45\textwidth}
        \begin{figure}[H]
            \centering
            \begin{tikzpicture}[node distance=3cm, auto, thick]
                \node[state] (S) {\emoji{sun}};
                \node[state] (R) [right of=S] {\emoji{umbrella}};

                \path[->]
                (S) edge[bend left] node {0.2} (R)
                (R) edge[bend left] node {0.4} (S)
                (S) edge[loop above] node {0.8} (S)
                (R) edge[loop above] node {0.6} (R);
            \end{tikzpicture}
            % \caption*{Gerichteter Graph}
        \end{figure}
        \begin{itemize}
            \item[{\faCircle[regular]}]<2->  \textbf{Knoten} = Zustände
            \item[\faArrowRight]<3-> \textbf{Kanten} = Übergänge mit Wahrscheinlichkeiten
        \end{itemize}
        \only<4->{
            Summe der ausgehenden Kanten = 1!
        }

        \column{.6\textwidth}
        \only<5->{
            \begin{myexercise}
                \small
                Jedes Mal, wenn Sie Instagram (\textbf{I}) benutzen, benutzen Sie als nächste App zu \textbf{60\%} nochmals Instagram, und wechseln auf TikTok (\textbf{T}) mit \textbf{40\%}. Wenn Sie zuerst TikTok benutzen, wechseln Sie zu \textbf{10\%} zu Instagram und bleiben zu \textbf{90\%} auf TikTok. Erstellen Sie den \textbf{gerichteten Graphen} für die Zustände \textbf{I} und \textbf{T}.
            \end{myexercise}
        }
    \end{columns}
\end{frame}

\begin{frame}{Gerichteter Graph: Social-Media-App-Nutzung}
    \begin{figure}[H]
        \centering
        \begin{tikzpicture}[node distance=4cm, auto, thick]
            \node[state] (I) {I};
            \node[state] (T) [right of=I] {T};

            \path[->]
            (I) edge[bend left] node[above] {0.4} (T)
            (I) edge[loop above] node {0.6} (I)
            (T) edge[bend left] node[below] {0.1} (I)
            (T) edge[loop above] node {0.9} (T);
        \end{tikzpicture}
        \caption*{Gerichteter Graph für \textbf{Instagram (I)} und \textbf{TikTok (T)}}
    \end{figure}
\end{frame}

\begin{frame}{Stationäre Verteilung}
    \begin{mydefinition}[Stationäre Verteilung]
    Eine Wahrscheinlichkeitsverteilung $\pi = (\pi_1, \pi_2, \ldots)$ heisst \textbf{stationäre Verteilung} einer Markov-Kette mit Übergangsmatrix $P$, wenn gilt:
    \[
        \pi = \pi \cdot P
    \]
    \end{mydefinition}

    \begin{itemize}[<+->]
        \item[$\Rightarrow$] Langfristige, stabile Verteilung der Wahrscheinlichkeiten
        \item[$\Rightarrow$] Bleibt nach Multiplikation mit $P$ unverändert
    \end{itemize}
\end{frame}

\begin{frame}{Berechnung der stationären Verteilung (Wetter-Beispiel)}
    \begin{columns}
        \column{0.6\textwidth}
        \begin{enumerate}[<+->]
            \item $\pi = \pi \cdot P$ aufstellen: \[(\pi_{\text{\emoji{sun}}}, \pi_{\text{\emoji{umbrella}}}) = (\pi_{\text{\emoji{sun}}}, \pi_{\text{\emoji{umbrella}}}) \cdot \begin{pmatrix} 0.8 & 0.2 \\ 0.4 & 0.6 \end{pmatrix}\]
            \item Gleichungen aufstellen:
                  \[
                      \begin{aligned}
                          \pi_{\text{\emoji{sun}}}      & = \pi_{\text{\emoji{sun}}} \cdot 0.8 + \pi_{\text{\emoji{umbrella}}} \cdot 0.4 \\
                          \pi_{\text{\emoji{umbrella}}} & = \pi_{\text{\emoji{sun}}} \cdot 0.2 + \pi_{\text{\emoji{umbrella}}} \cdot 0.6
                      \end{aligned}
                  \]
            \item Umformen:
                  \[
                      \begin{aligned}
                          \pi_{\text{\emoji{sun}}} \cdot 0.2   & = \pi_{\text{\emoji{umbrella}}} \cdot 0.4 \\
                          \Rightarrow \pi_{\text{\emoji{sun}}} & = 2 \cdot \pi_{\text{\emoji{umbrella}}}
                      \end{aligned}
                  \]
        \end{enumerate}
        \column{.01\textwidth}
        \begin{center}
            \vrule width 1pt height 6cm
        \end{center}
        \column{0.5\textwidth}
        \begin{enumerate}[<+->]
            \setcounter{enumi}{3}
            \item Normierung:
                  \[
                      \begin{aligned}
                          \pi_{\text{\emoji{sun}}} + \pi_{\text{\emoji{umbrella}}}              & = 1           \\
                          2 \cdot \pi_{\text{\emoji{umbrella}}} + \pi_{\text{\emoji{umbrella}}} & = 1           \\
                          3 \cdot \pi_{\text{\emoji{umbrella}}}                                 & = 1           \\
                          \pi_{\text{\emoji{umbrella}}}                                         & = \frac{1}{3}
                      \end{aligned}
                  \]
            \item Ergebnis:
                  \[
                      \begin{aligned}
                          \pi_{\text{\emoji{sun}}}      & = \frac{2}{3} \\
                          \pi_{\text{\emoji{umbrella}}} & = \frac{1}{3}
                      \end{aligned}
                  \]
            \item[$\Rightarrow$] \textbf{Langfristig ist es an 2/3 der Tage sonnig und an 1/3 der Tage regnerisch.}
        \end{enumerate}

    \end{columns}
\end{frame}